\begin{figure}[H]
\centering
\begin{tikzpicture}[x=1.65cm,y=1.1cm,>=Stealth,font=\small]
 \draw[->] (0,0)--(4.65,0) node[right] {$t$};
 \draw[->] (0,0)--(0,3.65) node[above] {modulus};
 \foreach \x in {1,2,3,4} {
   \draw (\x,.04)--(\x,-.04) node[below] {$\x$};}
 \foreach \y in {1,2,3} {
   \draw (.04,\y)--(-.04,\y) node[left] {$\y$};}
 \draw[gray!50,dashed] (2,0)--(2,1);
 \draw[blue!70!black,very thick] (0,0)--(2,1)--(4.3,3.3);
 \draw[orange!80!black,thick,dashed] (0,0)--(2,0)--(4.3,2.3);
 \node[blue!70!black,anchor=west] at (.2,3.25)
   {$\whdelta_E(t)=\max\{t/2,t-1\}$};
 \node[orange!80!black,anchor=west] at (.2,2.7)
   {$\bardelta_E(t)=\max\{0,t-2\}$};
\end{tikzpicture}
\caption{The exact moduli of the original norm on $E$.
The averaged midpoint modulus is positive for every positive radius,
whereas the one-sided modulus vanishes through radius $2$.
Theorem~\ref{thm:exact} gives both equalities at each unit center.
These curves do not describe other equivalent norms.}
\label{fig:moduli}
\end{figure}
